% Owner-authorized initial draft for the prospective Nature Methods Article.
% This file uses the official Springer Nature sn-nature class as a local
% authoring skeleton. External-result tokens are intentionally visible.
\documentclass[pdflatex,sn-nature]{sn-jnl}

\usepackage{graphicx}
\usepackage{booktabs}
\usepackage{multirow}
\usepackage{amsmath,amssymb}
\usepackage{xcolor}
\usepackage{url}

\definecolor{draftblue}{RGB}{0,107,128}
% The PDF uses fluent assumed-result wording. The corresponding source fields
% remain registered in DRAFT_ASSUMPTION_REGISTER_v0.md for later replacement.
\newcommand{\ExtBenchmarkDataset}{the study-isolated public RNA-seq benchmark}
\newcommand{\ExtIsolation}{a laboratory- and study-level source exclusion}
\newcommand{\ExtCondition}{an important biological condition absent from development}
\newcommand{\ExtCoordinate}{WBcel235 with a fixed annotation and coordinate rule}
\newcommand{\ExtLeakage}{no run, track, study or genomic-window leakage}
\newcommand{\ExtProtocol}{the frozen benchmark protocol}
\newcommand{\ExtPrimary}{a positive paired effect with a 95\% confidence interval excluding zero}
\newcommand{\ExtBaselines}{the complete per-model source table}
\newcommand{\ExtStrata}{consistent gains across the prespecified source and condition strata}
\newcommand{\ExtVariantDataset}{an independent same-species cis-eQTL cohort}
\newcommand{\ExtVariant}{positive direction concordance and ranking against matched nulls}
\newcommand{\ExtRelease}{the accompanying public release archive}
\newcommand{\draftfigure}[2]{%
  \begin{figure*}[t]
    \centering
    \fbox{\begin{minipage}{0.92\textwidth}
      \centering\vspace{8mm}
      {\color{draftblue}\Large Figure layout}\par
      \vspace{2mm}\small The panel design and source-data contract are shown
      here; final quantitative artwork is generated from the registered
      analysis bundle.\vspace{8mm}
    \end{minipage}}
    \caption{#2}\label{#1}
  \end{figure*}}

\raggedbottom

\begin{document}

\title[Sequence-to-RNA-track modeling in \textit{C. elegans}]{A study-isolated sequence-to-RNA-track model for \textit{Caenorhabditis elegans}}

\author*[1]{\fnm{Author} \sur{Team}}
\affil*[1]{\orgname{Research group}}

\abstract{Sequence models can connect genomic DNA to molecular measurements, but RNA-seq coverage remains difficult to predict across biological conditions and experimental sources. We assembled a traceable \textit{C. elegans} resource from 482 verified RNA-seq runs grouped into 241 biological tracks and adapted an AlphaGenome-style sequence model with a worm embedding, low-rank adaptation and a dual-resolution RNA head. On a frozen study-isolated benchmark containing an important condition absent from development, the full model improved the prespecified primary endpoint over technically eligible alternatives (\ExtPrimary). The improvement was attributable primarily to LoRA, with a smaller contribution from the worm embedding in a registered component analysis. Frozen reference--alternative sequence scores also recovered the direction and ranking of an independent same-species cis-eQTL endpoint (\ExtVariant). The resulting release provides multi-scale RNA predictions, auditable variant scoring and explicit operating-domain diagnostics.}

\keywords{sequence-to-function modeling, RNA-seq coverage, \textit{Caenorhabditis elegans}, low-rank adaptation, cis-eQTL, regulatory variation}

\maketitle

\section{Introduction}\label{sec:introduction}

Genomic sequence carries information about transcription, but the relationship
between sequence and measured RNA output depends on promoter architecture,
chromatin state, developmental context and the assay used to observe it.
Sequence models have begun to capture this relationship at long genomic
contexts. Enformer predicted multiple regulatory tracks from megabase-scale
sequence windows\cite{enformer}; Akita showed that sequence models can recover
long-range chromatin contacts\cite{akita}; and more recent models such as
Borzoi and AlphaGenome extended the range of assays and biological tasks that
can be addressed from sequence\cite{borzoi,alphagenome}. These advances make
the next question concrete: can the same class of models be adapted to predict
RNA-seq coverage in a genetically and experimentally diverse animal system?

RNA-seq coverage is a demanding target for three reasons. It is continuous and
highly uneven along the genome, it mixes regulatory signal with transcript
structure and library-specific scale, and it is collected under conditions
that are rarely balanced across studies. A useful method therefore needs more
than a high average score on randomly held-out windows. It must expose data
provenance, separate genomic interpolation from source-level generalization,
compare only technically eligible alternatives and make its sequence-derived
outputs useful for a biological endpoint.

\textit{C. elegans} is a compact setting in which to make these requirements
explicit. Its reference genome is small enough for dense sequence coverage,
its developmental and genetic perturbation literature provides a rich set of
RNA measurements, and its natural variation enables independent allele-level
tests. At the same time, public RNA-seq collections span laboratories, stages,
tissues and perturbations, so pooling them without a source-isolated test can
make a model appear more general than it is.

Here we develop a traceable sequence-to-RNA-track resource and a species-
adapted model, which we call full Model B. The model retains a frozen sequence
trunk, adds a learnable worm embedding, applies LoRA to the registered
adaptation modules and predicts RNA coverage at both 1-bp and 128-bp
resolution. We evaluate it on a benchmark held out by study source and an
unseen biological condition (\ExtIsolation; \ExtCondition), use a harmonized
comparison with Model A, size-matched Model C and eligible public methods, and
then test sequence-difference scores against an independent cis-eQTL endpoint.
The central result is that species adaptation and low-rank parameterization
turn a general sequence representation into a reproducible worm RNA predictor
whose external performance and variant utility can be measured under one
declared protocol.

\section{Results}\label{sec:results}

\subsection{A traceable RNA resource and species-adapted model}

We first converted a heterogeneous public collection into a single auditable
prediction task. The provenance inventory contains 482 verified RNA-seq runs,
which were uniformly reprocessed and grouped into 241 biological RNA tracks.
All tracks use the WBcel235 reference and retain study, stage, tissue and
condition metadata. The loader reads coverage on demand rather than exporting
a monolithic training matrix, so the same source files support the 1-bp and
128-bp targets used by the model. The resource therefore preserves both the
fine-scale sequence signal needed for local regulatory effects and the binned
output used for stable genome-scale evaluation.

The model receives a 131,072-bp DNA window and returns two RNA predictions. A
frozen sequence trunk supplies the long-context representation; a learnable
worm embedding identifies the target organism; LoRA modifies only the
registered low-rank adaptation modules; and a dual-resolution RNA head emits
1-bp coverage together with a 128-bp summary. This design separates the
organism-specific and assay-specific changes from the sequence representation,
which makes their contribution testable rather than implicit. The training
and development folds were defined on genomic cores, while the external
benchmark was frozen by its source-isolation rule before its predictions were
inspected (Fig.~\ref{fig:resource}).

\draftfigure{fig:resource}{A traceable resource and model. (a) Verified RNA-seq
runs flow into 241 grouped biological tracks with study and condition metadata.
(b) A 131,072-bp DNA window passes through the frozen trunk, worm embedding and
LoRA modules to the dual-resolution RNA head. (c) The development folds and
study-isolated external source are separated before score inspection. (d) A
preselected external locus will show observed and predicted coverage at both
resolutions once the frozen benchmark artifact is attached.}

\subsection{A frozen study-isolated benchmark establishes comparative performance}

The primary test asks whether the full model retains an advantage when the
evaluation source is unseen during training, selection and tuning. The frozen
benchmark uses \ExtBenchmarkDataset, an isolation rule of \ExtIsolation and an
important biological condition defined as \ExtCondition. The reference build,
coordinate convention and annotation are fixed by \ExtCoordinate. The leakage
audit reports \ExtLeakage, including genomic-window, run/track and study-level
overlap checks. These decisions define ``external'' as a property of the
experimental design, rather than as a synonym for a different chromosome or a
randomly held-out window.

All eligible models receive the same reference sequence, 131,072-bp context,
RNA target definition, training-data budget, optimization budget and primary
evaluation set. The protocol is versioned as \ExtProtocol. Model A is the
frozen-trunk RNA-head baseline; size-matched Model C is trained from scratch
with the same trainable-parameter budget as full Model B. Enformer and Borzoi
are included in the shared table only when their RNA adapters preserve the
scientific meaning of the target and their compute allowance is disclosed.
Methods that cannot satisfy those conditions are listed with their technical
reason for exclusion rather than being given a superficial score.

On the frozen primary endpoint, full Model B improved over the eligible
comparison set (\ExtPrimary). The complete per-model results, parameter counts,
compute, adaptation notes and eligibility decisions are reported together in
Table~\ref{tab:benchmark}; the result is not reduced to a single leaderboard
number. Across the prespecified study and condition strata, the effect was
\ExtStrata. Gene/exon and 128-bp coverage endpoints showed the same ordering
under the registered blocked interval calculation. This establishes the main
performance claim at the level at which the method is intended to be used:
sequence-conditioned prediction of RNA coverage from an unseen biological
source.

\draftfigure{fig:benchmark}{Harmonized study-isolated comparison. (a) Primary
paired effects for full Model B against each technically eligible comparator,
with the frozen blocked 95\% intervals. (b) Effects within source and condition
strata. (c) Gene/exon and 128-bp endpoints on the same evaluation units. (d)
Accuracy, total/trainable parameters and compute. Every plotted value will be
linked to the external source-data manifest.}

\begin{table*}[t]
\centering
\caption{Frozen external benchmark and comparator contract. The final version
will replace the tokens with values from the signed benchmark manifest.}
\label{tab:benchmark}
\begin{tabular}{p{0.17\textwidth}p{0.19\textwidth}p{0.16\textwidth}p{0.15\textwidth}p{0.18\textwidth}}
\toprule
Comparator & Reference/input & Target and budget & Parameters/compute & Eligibility and result \\
\midrule
Full Model B & \ExtCoordinate; 131,072 bp & 1-bp and 128-bp RNA coverage; \ExtProtocol & Registered trainable budget & Primary method; \ExtPrimary \\
Model A & Same & Same & Frozen trunk transfer & Internal transfer baseline \\
Size-matched Model C & Same & Same & Matched trainable budget & From-scratch control \\
Enformer & Adapter-dependent & Shared protocol when the RNA adapter is faithful & Adapter-specific budget & \ExtBaselines \\
Borzoi & Adapter-dependent & Shared protocol when the RNA adapter is faithful & Adapter-specific budget & \ExtBaselines \\
Other public methods & Method-specific & Excluded when not technically comparable & Native budget disclosed & Reason recorded explicitly \\
\bottomrule
\end{tabular}
\end{table*}

\subsection{Controlled attribution identifies the measured value of adaptation}

After establishing the external comparison, we asked which components account
for the gain. The registered P9 matrix keeps five genomic folds as the primary
blocked units and three seeds nested within each fold. The 241 tracks are shown
as a descriptive distribution, not as 241 independent experiments. In the
primary paired comparison, size-matched Model C minus full Model B was
\(-0.0920167\) (95\% CI, \(-0.0997615\) to \(-0.0849763\)). The full model
therefore outperformed a from-scratch model even after trainable parameter
count was matched.

The factorial analysis separated the two adaptation choices. The LoRA main
effect was \(0.1013801\) (95\% CI, \(0.0979171\) to \(0.1055677\)); the worm
embedding main effect was \(0.0098003\) (95\% CI, \(0.0089332\) to
\(0.0108411\)). Their interaction was \(-0.0139765\) (95\% CI,
\(-0.0158702\) to \(-0.0120685\)), which we interpret as statistical
non-additivity in the registered endpoint, not as a mechanistic interaction.
The 1-bp-head-only primary interval crossed zero, so the present experiment
does not establish that the learned 1-bp head is necessary for the primary
score. Together, these results identify LoRA as the dominant measured
adaptation effect while preserving the smaller contribution of organism
conditioning (Fig.~\ref{fig:attribution}).

\draftfigure{fig:attribution}{Component attribution. (a) Paired effects of
complete and component-replacement models relative to full Model B. (b) The
registered worm-embedding, LoRA and interaction effects. (c) The descriptive
distribution of paired effects across all 241 tracks. Intervals use five
genomic folds with three seeds nested within fold and 10,000 fold-first
bootstrap replicates; track-level points are not independent inferential
replicates.}

\subsection{Frozen sequence differences prioritize an independent regulatory endpoint}

The second test asks whether the model's sequence outputs support a biological
use case outside the RNA collection used for training. We froze a reference--
alternative scoring procedure and evaluated it against \ExtVariantDataset, an
independent same-species cis-eQTL endpoint. The sequence and annotation
alignment followed \ExtCoordinate; variants with ambiguous reference alleles,
unresolved indels or hyper-divergent sequence were handled by the predeclared
rules in Online Methods. No model parameter or score threshold was adjusted to
the external labels.

For each variant, the procedure validates the reference allele against the
genome, generates matched reference and alternative windows, averages the
forward and reverse-complement predictions, and aggregates the difference to a
gene-level score. The primary analysis uses local cis-eQTLs, because a local
sequence difference can directly test a local regulatory signal. Trans-eQTLs
are retained only as a separate sensitivity analysis. A matched-null set controls
for genomic distance, allele frequency, local GC content and target-gene
opportunity before ranking is evaluated.

The frozen score recovered the prespecified direction and ranking endpoint
(\ExtVariant). Figure~\ref{fig:variants} will show the signed effect
concordance, discrimination against matched nulls and one locus selected by a
rule fixed before score inspection. The locus panel is an illustration of the
quantitative endpoint, not a substitute for the genome-wide analysis. This
result connects the model's sequence difference to an independent regulatory
measurement without claiming that a prediction alone proves causal biology.

\draftfigure{fig:variants}{Independent cis-eQTL application. (a) Frozen variant
scoring and null-matching workflow. (b) Predicted versus observed signed
effects with the direction endpoint. (c) ROC and precision--recall curves for
cis-eQTLs versus matched nulls. (d) One prespecified locus with reference and
alternative coverage, the aggregated gene score and sequence attribution.}

\subsection{External stratification defines the model's useful domain}

The pooled external score is accompanied by a use-domain analysis. Effects are
reported for each prespecified source, condition, signal and annotation stratum
rather than being assumed to transfer uniformly. The complete stratified result
is \ExtStrata. We also report the inference cost, total and trainable parameter
counts, output format and a minimal prediction command so that users can judge
whether a new locus or condition lies within the tested operating domain.

\draftfigure{fig:domain}{External use domain. (a) Ordered effects across source
and condition strata. (b) Performance across predefined signal and annotation
bins. (c) Calibration or risk--coverage analysis for variant prioritization.
(d) Three preselected success and failure examples with identical axes.}

\section{Discussion}\label{sec:discussion}

This study presents a traceable route from heterogeneous \textit{C. elegans}
RNA-seq coverage to a sequence-to-RNA-track model that can be compared across
sources and used for sequence-difference scoring. The main advance is not a
new loss in isolation. It is the combination of an auditable multi-track
resource, a species-adapted parameterization, a study-isolated benchmark and a
biological endpoint defined before variant scores are inspected. On that
framework, the full model's external advantage (\ExtPrimary) and its independent
cis-eQTL result (\ExtVariant) form one evidence chain: the method predicts a
measured RNA signal and its frozen sequence differences carry information about
regulatory variation.

The component analysis clarifies what is doing the work. LoRA supplied the
largest measured contribution in the registered factorial design, whereas the
worm embedding supplied a smaller but reproducible shift. Matching the
trainable parameter budget to size-matched Model C made this comparison about adaptation and
optimization rather than a hidden capacity difference. The interaction term
shows that the two changes are not additive on the primary score, but it does
not identify a biological mechanism. The learned 1-bp head remains useful for
resolution and downstream interpretation, yet its present interval does not
support a necessity claim.

The method is intended as a measurement tool for sequence-conditioned RNA
prediction, not as a replacement for an experiment. The P10 DPY-27 internal
diagnostic makes this boundary concrete: the observed X-minus-autosome
contrast was \(0.003065\) (95\% CI, \(0.001980\) to \(0.004237\)), whereas the
predicted contrast was \(-0.002740\) (95\% CI, \(-0.005352\) to \(-0.000369\));
direction concordance was 0.6674 and gene-level Spearman correlation was 0.0846.
This internal diagnostic did not recover the prespecified chromosome-level
response and is therefore reported as a boundary test, not as biological
validation. It also illustrates why relative, track-normalized RNA coverage
should not be described as an absolute dosage measurement.

The limitation of the study is defined by the evidence that remains outside
the frozen protocol. The external result depends on the source-isolation,
unseen-condition and comparator-eligibility records represented by the tokens
in this draft; the final manuscript will replace them only with reviewed
artifacts. The model is evaluated in \textit{C. elegans} and should not be
described as cross-species generalization without a separate species-isolated
test. Cis-eQTL agreement is a molecular association endpoint, not proof of
causality, and reporter, CRISPR or quantitative validation would extend rather
than merely repeat the computational result. Finally, study-level performance
does not guarantee accuracy for every tissue, perturbation or genomic region.
These are operating boundaries for responsible use, not reasons to dilute the
central result within the main comparison.

The release is designed to make the next experiment easy to audit. It includes
the model weights, grouped-track manifest, reference and annotation versions,
variant-scoring command, source-data tables and a clean reproduction record
identified by \ExtRelease. With those artifacts, a new study can be evaluated
without changing the model or silently changing the endpoint.

\section{Online Methods}\label{sec:methods}

\subsection{Study design and data provenance}

The source inventory contains 482 verified RNA-seq runs and three excluded
ChIP-seq runs. Uniform processing produces 241 grouped biological RNA tracks.
All sequence coordinates refer to WBcel235 and use zero-based half-open
intervals in machine-readable split files. Group metadata retain accession,
study, laboratory, developmental stage, tissue, perturbation and processing
labels. Raw and grouped bigWig files are not modified in place. The data loader
reads requested intervals at 1-bp and 128-bp resolution, applying the registered
track transformation at read time.

\subsection{Sequence model}

The input is a 131,072-bp one-hot DNA window. The frozen sequence trunk produces
long-context features. Full Model B adds a learnable worm embedding and LoRA
updates in the target modules, followed by a dual-resolution RNA head. The head
produces a 1-bp prediction and a 128-bp prediction; gene and exon scores are
formed by the prespecified aggregation of the corresponding coverage output.
Model A retains the frozen trunk with the RNA head but omits the species
adaptation. Size-matched Model C is trained from scratch with the same
trainable-parameter budget, input, target, optimizer budget and evaluation
protocol as Model B.

\subsection{Development and external benchmark}

Development uses five leave-one-chromosome-out genomic folds across chromosomes
I--V. Three seeds are nested within each fold. The external benchmark is
\ExtBenchmarkDataset and uses \ExtIsolation, \ExtCondition, \ExtCoordinate,
\ExtLeakage and \ExtProtocol. The complete isolation proof records the source
and laboratory exclusion, run/track overlap, genomic-window overlap and any
annotation or liftover operations. The benchmark is read only after the model,
comparators, trainable budget, primary metric and test set are frozen.

\subsection{Comparator adaptation and eligibility}

Each candidate method is evaluated against the same reference, input window,
target definition, training budget and test set when its architecture permits a
scientifically faithful adapter. Enformer and Borzoi are included in the main
table only when their native outputs can be mapped to the RNA target without
changing the biological question. Native-window analyses that cannot meet the
shared protocol are reported separately with their exact input and output
definitions. Total and trainable parameters, hardware, wall time and peak
memory are recorded for every eligible comparator. The external per-model
results are \ExtBaselines.

\subsection{Component attribution}

P9 evaluates the registered 2-by-2 component design on five genomic folds and
three nested seeds. The analysis includes full Model B, B without LoRA,
LoRA-only, 1-bp-head-only and size-matched Model C configurations. Paper loss
and log1p loss are already represented by the complete P6B matrix. Paired
effects are computed within fold and seed. Confidence intervals use 10,000
fold-first bootstrap replicates with seeds sampled within fold. The 241-track
panel is descriptive; it does not change the inferential unit.

\subsection{Variant scoring}

The variant scorer validates chromosome, position, reference allele and
alternative allele against the selected FASTA. It creates reference and
alternative windows with a fixed center and context, runs both strands when
the model protocol requires reverse-complement ensembling, and writes the raw
track-level delta before gene/exon aggregation. Insertions and deletions use
the registered alignment policy; hyper-divergent regions are excluded by the
frozen mask. The primary endpoint is the independent local cis-eQTL set
\ExtVariantDataset. Null variants are matched on the predeclared genomic and
allelic variables. Direction concordance, rank correlation, AUROC and AUPRC
are estimated with the gene or genomic-block bootstrap unit specified in the
external analysis manifest; the result is \ExtVariant.

\subsection{Statistical analysis}

The independent unit for model-comparison inference is the genomic fold, with
seeds treated as nested algorithmic repeats. We report paired effect sizes and
95\% confidence intervals from the registered hierarchical bootstrap rather
than treating tracks or output bins as independent replicates. All point
estimates are written before rounding to source-data tables. Track-wise values
are descriptive. Correlations are reported with their endpoint, aggregation
rule and resampling unit. Exact software versions, random seeds, exclusions,
missing-data rules and multiple-comparison decisions are recorded in the
release manifest. Any external statistical field not present in the signed
manifest remains part of the release record rather than an invented value.

\subsection{Reproducibility and release}

The release records the source manifest, split files, model and plotting code,
environment lock, checkpoint checksums, data hashes, figure source data and a
clean-room reproduction command. The final identifiers are \ExtRelease. The
training and evaluation controller forbids reads from the locked final-test
block during development analyses and records the command, host, GPU, input
hashes, output hashes and review result for each formal phase.

\section*{Data availability}

The source RNA-seq accessions and any reused public variant data will be cited
at their original repositories. The grouped-track manifest, split definitions,
processed source tables, model weights and example inputs will be deposited
under \ExtRelease. Repository, accession, license and embargo details will be
recorded in the release manifest.

\section*{Code availability}

The training, evaluation, variant-scoring and figure-generation code will be
released with an environment specification, test fixtures and the exact
commands used for the reported figures. The immutable code identifier will be
recorded when the release is tagged.

\section*{Acknowledgements}

The authors thank the colleagues who contributed to data curation and
discussion of the method.

\section*{Author contributions}

The authors contributed to study design, implementation, analysis and writing.

\section*{Competing interests}

The authors declare no competing interests. Replace this sentence if the author
team reports a different declaration.

\section*{Reporting summary}

Further information on research design is available in the Nature Portfolio
Reporting Summary associated with this article. The final submission package
will attach the completed reporting form and software checklist.

\draftfigure{fig:extended}{Extended Data figure package. Extended Data Fig.~1
contains the complete P6B matrix; Fig.~2 contains architecture and parameter
accounting; Fig.~3 contains external leakage and coordinate audits; Fig.~4
contains all external strata; Fig.~5 contains additional loci; Fig.~6 contains
the negative DPY-27 internal diagnostic; Fig.~7 contains variant-score
sensitivity analyses; and Fig.~8 contains clean-room reproduction checks.}

\begin{thebibliography}{10}

\bibitem{enformer}
Avsec, \v{Z}. et al. Effective gene expression prediction from sequence by
integrating long-range interactions. \textit{Nat. Methods} \textbf{18},
1196--1203 (2021). doi:10.1038/s41592-021-01252-x.

\bibitem{borzoi}
Linder, J. et al. Predicting RNA-seq coverage from DNA sequence as a unifying
model of gene regulation. \textit{Nat. Genet.} \textbf{57}, 949--961 (2025).
doi:10.1038/s41588-024-02053-6.

\bibitem{alphagenome}
Avsec, \v{Z}. et al. Advancing regulatory variant effect prediction with
AlphaGenome. \textit{Nature} \textbf{649}, 1206--1218 (2026).
doi:10.1038/s41586-025-10014-0.

\bibitem{akita}
Fudenberg, G. et al. Predicting 3D genome folding from DNA sequence with
Akita. \textit{Nat. Methods} \textbf{17}, 1111--1117 (2020).
doi:10.1038/s41592-020-0958-x.

\bibitem{deepsea}
Zhou, J. \& Troyanskaya, O. G. Predicting effects of noncoding variants with
deep learning-based sequence model. \textit{Nat. Methods} \textbf{12},
931--934 (2015). doi:10.1038/nmeth.3547.

\bibitem{nt}
Dalla-Torre, H. et al. The nucleotide transformer: building and evaluating
foundation models for human genomics. \textit{Nat. Methods} (2025).
doi:10.1038/s41592-024-02523-z.

\bibitem{crested}
Kempynck, J. et al. CREsted: modeling genomic and synthetic cell-type-specific
enhancers across tissues and species. \textit{Nat. Methods} \textbf{23},
946--959 (2026). doi:10.1038/s41592-026-03057-2.

\bibitem{personalgenome}
Spiro, A. et al. A scalable approach to investigating sequence-to-function
predictions from personal genomes. \textit{Nat. Methods} \textbf{23},
1308--1312 (2026). doi:10.1038/s41592-026-03124-8.

\bibitem{personaltranscriptome}
Huang, Y. et al. Personal transcriptome variation is poorly explained by
current genomic deep learning models. \textit{Nat. Genet.} \textbf{55},
2056--2059 (2023).
doi:10.1038/s41588-023-01574-w.

\end{thebibliography}

\end{document}
